% Generated from validated Article Draft Result. Do not hand-edit; revise the structured draft instead.
\section{Agent Safety \& Security — 「使える」と「許される」を分ける}
\label{pkg:agent-safety-security-june}
\noindent\textbf{agentがtool、memory、外部serviceへ接続するほど、securityはprompt injectionだけでは説明できなくなる。6月の研究はcapability exposure、per-call authorization、MCP tool poisoning、persistent-agent attack surface、adaptive evaluationを別々の問題として切り分けた。安全性評価も、固定attack suiteから実行系全体へ広がっている。}\autocite{src-02da0db00c2d,src-6da0c5a88fa4,src-91b06074799d,src-c000018ba1f0,src-cd4b60edfee7}

\subsection{Capability gateはAuthorizationではない}

Capability Gates Are Not Authorizationの著者らは、agentにtool capabilityを見せる／隠す制御と、実際の各tool callで値やscopeを検証するauthorizationは別物だと整理した。一般的なagent frameworkを対象にconfused-deputy riskを示し、fail-closedなper-call controlの必要性を論じている。\autocite{src-02da0db00c2d}


\subsection{防御を知ったattackで評価する}

prompt injection防御の評価では、固定attackを同じ形で当て続けるだけでは防御の実力を測りにくい。Adaptive Evaluationの著者らは、defense-awareなattackへ適応させる評価を主張し、Progent/AgentDojoを用いた検証を行った。防御を公開した後の攻撃者適応まで評価modelへ入れる考え方である。\autocite{src-cd4b60edfee7}


\subsection{persistent agentはcomputer systemとして攻撃面が増える}

SafeClawArenaの著者らはClaw-likeなpersistent agentをcomputer systemとして捉え、単一promptだけでなく複数componentにまたがるattack surfaceを評価対象にした。長時間稼働し、toolやstateを持つagentでは、model inputだけを守ってもsystem全体のsecurityを説明できないという問題設定である。\autocite{src-6da0c5a88fa4}


\subsection{MCPではtool集合そのものがtrust boundaryになる}

ShareLockの著者らはMCP環境で、複数toolに分散したpoisoning情報を組み合わせるthreshold-style attackを提示した。個々のtoolを単独で検査して安全に見えても、agentが複数toolを横断して情報を統合した時にattackが成立し得るというthreat assumptionが重要である。\autocite{src-91b06074799d}


\subsection{本番前に、過去の会話を候補modelへ再生する}

OpenAIはDeployment Simulationとして、過去のconversationをcandidate modelとsimulated tool環境へreplayし、pre-deploymentで振る舞いを観察する方法を紹介した。これは運用前評価の一つの手法としてのfirst-party報告であり、すべてのfailureを捕捉する保証と同義ではない。\autocite{src-c000018ba1f0}


\begin{center}
\small
\begin{tabularx}{\columnwidth}{>{\raggedright\arraybackslash}X|>{\raggedright\arraybackslash}X|>{\raggedright\arraybackslash}X}
\toprule
\textbf{問題} & \textbf{主な境界} & \textbf{代表的な6月材料} \\
\midrule
Confused deputy & toolを使えること vs 呼出しを許可すること & Capability Gates \\
Prompt injection & 固定attack vs defense-aware attack & Adaptive Evaluation \\
Persistent agent & model input vs system全体のattack surface & SafeClawArena \\
MCP poisoning & 単一tool検査 vs 複数toolの合成 & ShareLock \\
Pre-deployment safety & offline test vs 実際の利用履歴に近いsimulation & Deployment Simulation \\
\bottomrule
\end{tabularx}
\autocite{src-02da0db00c2d,src-6da0c5a88fa4,src-91b06074799d,src-c000018ba1f0,src-cd4b60edfee7}
\end{center}

これらを一括して「agentは危険」と結論づけるのは粗すぎる。権限設計、入力汚染、tool間の合成、persistent state、評価方法はfailure modeが異なり、必要なcontrolも異なる。共通するのは、model outputを信用するか否かではなく、実行前後に独立したpolicy boundaryを置く発想である。\autocite{src-02da0db00c2d,src-6da0c5a88fa4,src-91b06074799d,src-cd4b60edfee7}


\begin{claimboundary}
4本のsecurity paperで示されたattack、benchmark、防御効果はそれぞれ著者のthreat modelとevaluation環境に依存する。別frameworkや実運用へそのまま成功率・防御率を一般化せず、何を攻撃者へ許し、どの境界を検証したかを先に確認する。\autocite{src-02da0db00c2d,src-6da0c5a88fa4,src-91b06074799d,src-cd4b60edfee7}
\end{claimboundary}

\begin{claimboundary}
Deployment Simulationの有効性はOpenAIによるfirst-party safety methodの報告である。実運用への近さは評価設計上の利点になり得るが、独立した再現や完全なrisk coverageを意味しない。\autocite{src-c000018ba1f0}
\end{claimboundary}
